
\chapter{Resultados}
\label{chap:resultados}

% Apresentar os resultados obtidos ao longo do desenvolvimento
% do DoseCerta, considerando os objetivos e requisitos definidos.

\section{Visão Geral dos Resultados}
\label{sec:visao-geral-resultados}

% Resumir o estágio alcançado pelo projeto:
% protótipo construído, módulos implementados e funções disponíveis.
% Distinguir claramente o que foi concluído do que ficou pendente.

\section{Resultados do Desenvolvimento Mecânico}
\label{sec:resultados-mecanicos}

% Apresentar o resultado físico do projeto mecânico:
% componentes fabricados, montagem, dimensões finais e imagens.
% Incluir somente características confirmadas no protótipo.

\section{Resultados do Desenvolvimento Eletrônico e Embarcado}
\label{sec:resultados-eletronicos-embarcados}

% Registrar os circuitos montados e as funções embarcadas
% que foram implementadas e verificadas.
% Incluir medições ou evidências disponíveis.

\section{Resultados do Backend e do Aplicativo}
\label{sec:resultados-backend-aplicativo}

% Apresentar as funcionalidades efetivamente implementadas.
% Incluir telas, respostas da API, registros ou outras evidências,
% conforme disponíveis.

\section{Resultados da Visão Computacional e do Agente de IA}
\label{sec:resultados-visao-ia}

% Apresentar resultados observados nesses módulos, se desenvolvidos.
% Informar métricas somente quando houver dados e procedimento
% de avaliação que as sustentem.

\section{Resultados dos Testes}
\label{sec:resultados-testes}

% Sintetizar os resultados dos testes descritos no Capítulo 14.
% Indicar quais critérios foram atendidos, parcialmente atendidos
% ou não atendidos, com referência às evidências.

% Exemplo de tabela de síntese:
% \begin{table}[htbp]
%     \centering
%     \caption{Síntese dos resultados dos testes}
%     \label{tab:sintese-resultados-testes}
%     \begin{tabular}{|l|p{5cm}|p{3cm}|}
%         \hline
%         Teste & Critério avaliado & Resultado obtido \\
%         \hline
%         % T01 & Descrever o critério & A preencher \\
%         \hline
%     \end{tabular}
% \end{table}

\section{Atendimento aos Requisitos}
\label{sec:atendimento-requisitos}

% Relacionar os requisitos do Capítulo 5 às evidências
% obtidas durante a implementação e os testes.
% Registrar requisitos atendidos, parcialmente atendidos
% ou não atendidos, conforme os critérios estabelecidos.

\section{Síntese dos Resultados}
\label{sec:sintese-resultados}

% Retomar os principais resultados concretos do projeto
% e indicar o estágio final alcançado pelo DoseCerta.
